\documentclass[12pt,aspectratio=169]{beamer}
\input{header_beamer}
\usetheme{Boadilla}
\usecolortheme{beaver}
\setbeamertemplate{page number in head/foot}{}
\setbeamertemplate{navigation symbols}{}


\title{ Lecture-23: Invariant Distribution}
\date{}

\begin{document}
\pagestyle{empty}
\maketitle

\begin{frame}{Invariant Distribution} 
Let $X: \Omega \to \sX^{\Z_+}$ be a time-homogeneous Markov chain with transition probability matrix $P: \sX\times\sX\to[0,1]$. 
\begin{Definition}<2->
A probability distribution $\pi \in \cM(\sX)$ is said to be %\textbf{stationary distribution} or 
\textbf{invariant distribution} for the Markov chain $X$ if it satisfies the \textbf{global balance equation} $\pi = \pi P$. 
\end{Definition}
\begin{Definition}<3->
When the initial distribution of a Markov chain is $\nu \in \cM(\sX)$, 
then the conditional probability is denoted by $P_\nu: \sF \to [0,1]$ defined by 
\EQ{
P_\nu(A) \triangleq \sum_{x\in\sX}\nu(x)P_x(A)\text{ for all events }A \in \sF.
} 
\end{Definition}
\end{frame}

\begin{frame}
\begin{Definition} 
For a Markov chain $X:\Omega\to\sX^{\Z_+}$, 
we denote the distribution of random variable $X_n:\Omega\to\sX$ by $\nu_n \in \cM(\sX)$ for all $n \in \Z_+$. 
That is, $\nu_n(x) \triangleq P_{\nu_0}\set{X_n = x}$ for all $x \in \sX$. 
\end{Definition}
\begin{block}{Reminder}<2->
We observe that $\nu_n(x) = \sum_{z\in\sX}\nu_0(z)(P^n)_{zx}$ for all $x \in \sX$. 
\end{block}
\end{frame}

\begin{frame}
\begin{block}{Reminder} 
Facts about the invariant distribution $\pi$. 
\begin{itemize}
\item The global balance equation $\pi = \pi P$ is a matrix equation; that is, we have a collection of $\abs{\sX}$ equations: $\pi_y = \sum_{x \in \sX}\pi_x p_{xy}$ for each $y \in \sX$.

\item<2-> The invariant distribution $\pi$ is a left eigenvector of the stochastic matrix $P$ with the largest eigenvalue $1$. The all-ones vector is the right eigenvector of this stochastic matrix $P$ for the eigenvalue $1$.

\item<3-> From the Chapman-Kolmogorov equation for the initial probability vector $\pi$, we have $\pi = \pi P^n$ for $n \in \N$. That is, if $\nu_0 = \pi$, then $\nu_n = \pi$ for all $n \in \Z_+$.

\item<4-> The resulting process with initial distribution $\pi$ is stationary, and hence has shift-invariant finite-dimensional distributions. For example, for any $k, n \in \Z_+$ and $x_0, \dots, x_n \in \sX$, we have
\EQ{
P_{\pi}\set{X_0 = x_0, \dots, X_n = x_n} = P_{\pi}\set{X_k = x_0, \dots, X_{k+n} = x_n} = \pi_{x_0}p_{x_0x_1} \dots p_{x_{n-1}x_n}.
}
\end{itemize}
\end{block}
\end{frame}


\begin{frame}
\begin{block}{Reminder} 
\begin{itemize}
\item For an irreducible Markov chain, if $\pi_x > 0$ for some $x \in \sX$, then the entire invariant vector $\pi$ is positive. 
To this end, we will show that $\pi_y > 0$ for all states $y \in \sX$. 
Let $y \in \sX$, then from the irreducibility of Markov chain, there exists an $m \in \Z_+$ such that $p_{xy}^{(m)} > 0$. 
Further, $\pi = \pi P^m$ and hence $\pi_y \ge \pi_x p_{xy}^{(m)} > 0$.  

\item<2-> Any scaled version of $\pi$ satisfies the global balance equation. 
Therefore, for any invariant vector $\alpha \in \sX^{\R_+}$ of a positive recurrent transition matrix $P$,  
the sum $\norm{\alpha}_1 = \sum_{x \in \sX}\alpha_x$ must be finite. 
We can normalize $\alpha$ and get an invariant probability measure $\pi = \frac{\alpha}{\norm{\alpha}_1}$. 
\end{itemize}
\end{block}
\end{frame}

\begin{frame}
\begin{block}{Theorem} 
An irreducible Markov chain with transition probability matrix $P:\sX\times\sX\to[0,1]$ is positive recurrent iff there exists a unique invariant probability measure $\pi \in \cM(\sX)$ that satisfies global balance equation $\pi = \pi P$ and $\pi_x = \frac{1}{\mu_{xx}} > 0$ for all $x \in \sX$. 
\end{block}

\begin{Proof}<2->
Consider an irreducible Markov chain $X:\Omega\to\sX^{\Z_+}$ with transition probability matrix $P$.  
%We will first show that positive recurrence of $X$ implies the existence of a positive invariant distribution $\pi$ and its uniqueness. 
%Then we will show that the existence of a unique positive invariant distribution $\pi$ implies positive recurrence of $X$.  
\begin{itemize}
\item[Implication:]<3->
For Markov chain $X$, let the initial state be $X_0 = x$. 
Recall that the number of visits to state $y \in \sX$ in the first $n$ steps of the Markov chain $X$ is denoted by $N_y(n) = \sum_{k=1}^n\SetIn{X_k = y}$.  
It follows that $\sum_{y \in \sX}N_y(n) = n$ for each $n \in \N$. 
Let $H_x \triangleq \tau_X^{\set{x},1}$ be the first recurrence time to state $x\in\sX$, then we have $N_x(H_x) = 1$ and $\sum_{y \in \sX}N_y(H_x) = H_x$.
\end{itemize}
\end{Proof}
\end{frame}

\begin{frame}
\begin{block}{Proof (Existence)}
Define a vector $v \in \R^\sX$ by $v_y \triangleq \E_x[N_y(H_x)]$ for each $y \in \sX$. Observe that $v_y \ge 0$ for each state $y \in \sX$.  
In particular, $v_x = 1$ and $\sum_{y \in \sX} v_y = \E_x H_x = \mu_{xx} < \infty$.

\begin{itemize}
\item<2-> We will show that the vector $v$ satisfies the global balance equation $v = vP$. To see this, we first define
\[
\lambda_{xy}^{(n)}  \triangleq P_x\set{X_n = y, n \le H_x}
\]
for all $n \in \N$ and states $x,y \in \sX$.  
We observe that $\lambda_{xy}^{(1)} = p_{xy}$ for each $y \in \sX$.

\item<3-> Next, by the monotone convergence theorem:
\begin{equation*}
\label{eqn:Eigenvector}
v_y = \E_x N_y(H_x) = \E_x\sum_{n \in \N} \SetIn{X_n = y, n \le H_x}
= \sum_{n \in \N} P_x\set{X_n = y, n \le H_x} = \sum_{n \in \N} \lambda_{xy}^{(n)}.
\end{equation*}
\end{itemize}
\end{block}
\end{frame}

\begin{frame}
\begin{block}{Proof continued (Existence)}
\begin{itemize}
\item For $n \ge 2$, partitioning the event $\set{X_n = y, n \le H_x}$ by the events $(\set{X_{n-1} = z} : z \in \sX \setminus \set{x})$, and using the countable additivity and the definition of conditional probability, we get:
\[
\lambda_{xy}^{(n)} 
= \sum_{z \neq x} P\left(\set{X_n = y} \given \set{X_{n-1} = z, n \le H_x, X_0 = x}\right) P_x\set{X_{n-1} = z, n \le H_x}.
\]

\item<2-> Since $H_x$ is adapted to the natural filtration $\sF_\bullet$ of the Markov chain $X$, we have:
\[
\set{n \le H_x, X_0 = x} = \set{X_0 = x} \cap \set{H_x > n-1}^c \in \sF_{n-1}.
\]
\end{itemize}
\end{block}
\end{frame}


\begin{frame}
\begin{block}{Proof continued (Existence)}
\begin{itemize}
\item <1-> Together with the Markov property of $X$ and the fact that 
\[
\set{X_{n-1}=z, n \le H_x} = \set{X_{n-1}=z, n-1 \le H_x},
\]
we obtain,
\begin{align*}
%\label{eqn:EigenvectorTerms}
\lambda_{xy}^{(n)} 
&= \sum_{z \neq x} P\left( \set{X_n = y} \given \set{X_{n-1} = z} \right) P_x\set{X_{n-1} = z, n-1 \le H_x} \\
&= \sum_{z \neq x} \lambda_{xz}^{(n-1)} p_{zy}.
\end{align*}

\item <2-> Substituting the expression for $\lambda_{xy}^{(n)}$ into the expression for $v_y = \sum_{n \in \N} \lambda_{xy}^{(n)}$ from~\eqref{eqn:Eigenvector}, and using the fact that $v_x = 1$, we obtain:
\[
v_y = p_{xy} + \sum_{n \ge 2} \sum_{z \neq x} \lambda_{xz}^{(n-1)} p_{zy}
= v_x p_{xy} + \sum_{z \neq x} v_z p_{z y}
= \sum_{x \in \sX} v_x p_{x y}.
\]
\end{itemize}
\end{block}
\end{frame}

\begin{frame}
\begin{block}{Proof continued (Existence)}
\begin{itemize}
\item Since $v$ has a finite sum, it follows that
\[
\pi \triangleq \frac{v}{\sum_{x \in \sX} v_x}
\]
is an invariant distribution for the Markov chain $X$ with transition matrix $P$.

In addition, we have
\[
\pi_x = \frac{v_x}{\sum_{y \in \sX} v_y} = \frac{1}{\mu_{xx}} > 0.
\]

\end{itemize}
\end{block}
\end{frame}



\begin{frame}
\begin{block}{Proof Continued (Uniqueness)}
\begin{itemize}
    \item Next, we show that this is a unique invariant measure independent of the initial state $x$, 
    and hence $\pi_y = \frac{1}{\mu_{yy}} > 0$ for all $y \in \sX$.  
    For uniqueness, we observe from the Chapman-Kolmogorov equations and the invariance of $\pi$ that
    \[
    \pi = \frac{1}{n} \pi (P + P^2 + \dots + P^n).
    \]
    
    \item<2-> Hence,
    \[
    \pi_y = \sum_{x \in \sX} \pi_x \cdot \frac{1}{n} \sum_{k=1}^n p_{xy}^{(k)}, \quad \text{for all } y \in \sX.
    \]
    Taking the limit as $n \to \infty$ on both sides and exchanging the limit and summation on the right-hand side (justified by the bounded convergence theorem, since $\pi$ is summable), we obtain
    \[
    \pi_y = \frac{1}{\mu_{yy}} \sum_{x \in \sX} \pi_x = \frac{1}{\mu_{yy}} > 0, \quad \text{for all } y \in \sX.
    \]
\end{itemize}
\end{block}
\end{frame}

\begin{frame}
\begin{Proof}[Proof (Converse)]
Let $\pi$ be the unique positive invariant distribution of Markov chain $X$, such that $\pi_y = \frac{1}{\mu_{yy}} > 0$ for all states $y \in \sX$. 
It follows that the Markov chain $X$ is positive recurrent.
\end{Proof}
\end{frame}


\begin{frame}
\begin{cor} 
An irreducible Markov chain on a finite state space $\sX$ has a unique positive stationary distribution $\pi$. 
\end{cor}
\begin{Definition}<2->
An irreducible, aperiodic, positive recurrent Markov chain is called \textbf{ergodic}. 
\end{Definition}
\end{frame}

\begin{frame}
\begin{block}{Reminder} 
Additional remarks about the stationary distribution $\pi$. 
\begin{itemize}
    \item For a Markov chain with multiple positive recurrent communicating classes $\cC_1, \dots, \cC_m$, one can find the positive equilibrium distribution for each class and extend it to the entire state space $\sX$, denoting it by $\pi_k$ for class $k \in [m]$.
    
    \item<2-> It is easy to check that any convex combination
    \[
    \pi = \sum_{k = 1}^m \alpha_k \pi_k
    \]
    satisfies the global balance equation $\pi = \pi P$, where $\alpha_k \ge 0$ for each $k \in [m]$ and $\sum_{k=1}^m \alpha_k = 1$.
    
    \item<3-> Hence, a Markov chain with multiple positive recurrent classes has a convex set of invariant probability measures, with the individual invariant distributions $\pi_k$ for each positive recurrent class $k \in [m]$ being the extreme points. 
\end{itemize}
\end{block}
\end{frame}


\begin{frame}
\begin{block}{Reminder}
\begin{itemize}
\item Let $\mu(0) = e_x$, that is let the initial state of the positive recurrent Markov chain be $X_0 = x$. Then, we know that %for an aperiodic Markov chain $X$, 
\EQ{
\pi_y = \frac{1}{\mu_{yy}} = \lim_{n \to \infty}\frac{1}{n}\sum_{k=1}^np_{xy}^{(k)} =  \lim_{n \to \infty}\frac{1}{n}\E_xN_y(n).
}
That is, $\pi_y$ is limiting average of number of visits to state $y \in \sX$. 

\item<2-> If a positive recurrent Markov chain is aperiodic, then limiting probability of being in a state $y$ is its invariant probability, 
that is $\pi_y = \lim_{n \to \infty}p_{xy}^{(n)}$. 
\end{itemize}
\end{block}
\end{frame}

\begin{frame}
\begin{block}{Theorem} 
For an ergodic Markov chain $X$ with invariant distribution $\pi$, and $n$th step distribution $\mu(n)$, we have $\lim_{n \to \infty}\mu(n) = \pi$ in the total variation distance. 
\end{block}

\begin{block}{Proof}<2->
Consider independent time homogeneous Markov chains $X: \Omega \to \sX^{\Z_+}$ and $Y: \Omega \to \sX^{\Z_+}$ each with transition matrix $P$. 
The initial state of Markov chain $X$ is assumed to be $X_0 = x$, whereas the Markov chain $Y$ is assumed to have an initial distribution $\pi$. 
It follows that $Y$ is a stationary process, while $X$ is not. 
In particular, 
\meq{2}{
&\mu_y(n) = P_x\set{X_n = y} = p_{xy}^{(n)}, && P_{\pi}\set{Y_n = y} = \pi_y.
} 
Let $\tau = \inf\{n \in \Z_+: X_n = Y_n\}$ be the first time that two Markov chains meet, called the \textit{coupling time}. \\
\end{block}
\end{frame}

\begin{frame}
\begin{block}{Proof continued}
\begin{itemize}
\item \textbf{Finiteness:} First, we show that the coupling time is almost surely finite. 
To this end, we define a a new Markov chain on state space $\sX \times \sX$ with transition probability matrix $Q$ such that $q((x,w), (y,z)) = p_{xy}p_{wz}$ for each pair of states $(x,w), (y,z) \in \sX \times \sX$. 
The $n$-step transition probabilities for this couples Markov chain are given by 
\EQ{
q^{(n)}((x,w), (y,z)) \triangleq p^{(n)}_{xy}p^{(n)}_{wz}.
} 
\item<2-> \textbf{Ergodicity:} Since the Markov chain $X$ with transition probability matrix $P$ is irreducible and aperiodic, for each $x,y,w,z \in \sX$ there exists an $n_0 \in \Z_+$ such that $q^{(n)}((x,w), (y,z)) = p^{(n)}_{xy}p^{(n)}_{wz} > 0$ for all $n \ge n_0$ from a previous Lemma on aperiodicity. Hence, the irreducibility and aperiodicity of this new \textbf{product} Markov chain follows. 
\end{itemize}
\end{block}
\end{frame}

\begin{frame}
\begin{block}{Proof continued}
\begin{itemize}
\item \textbf{Invariant:} It is easy to check that $\theta(x,w) = \pi_x\pi_{w}$ is the invariant distribution for this product Markov chain, since $\theta(x,w) > 0$ for each $(x,w) \in \sX \times \sX$, $\sum_{x,w \in \sX}\theta(x,w) = 1$, and for each $(y,z) \in \sX \times \sX$, we have 
\EQ{
\sum_{x,w \in \sX}\theta(x,w)q((x,w), (y,z)) = \sum_{x \in \sX}\pi_xp_{xy}\sum_{w \in \sX}\pi_{w}p_{wz} = \pi_y\pi_{z} = \theta(y,z). 
}
\item<2-> \textbf{Recurrence:} This implies that the product Markov chain is positive recurrent, and each state $(x,x) \in \sX \times \sX$ is reachable with unit probability from any initial state $(y,w) \in \sX \times \sX$.  
\item<2-> In particular, the coupling time is almost surely finite. 
\end{itemize}
\end{block}
\end{frame}

\begin{frame}
\begin{Proof}[Proof continued]
\begin{itemize}
\item \textbf{Coupled process:} Second, we show that from the coupling time onwards, the evolution of two Markov chains is identical in distribution.  
That is, for each $y \in \sX$ and $n \in \Z_+$, 
\EQ{
P_{X_{\tau}}\set{X_{n} = y, n \ge \tau} = P_{Y_{\tau}}\set{Y_{n} = y, n \ge \tau}. 
}
%This follows from the strong Markov property for the joint process where 
$\tau$ is stopping time for the joint process $((X_n, Y_n): n \in \Z_+)$ such that $X_\tau = Y_\tau$, 
and both marginals  have the identical transition matrix.\\ % and that $X_{\tau} = Y_{\tau}$. 

\item<2->
\textbf{Limit:} For any $y \in \sX$, we can write the difference as 
\EQ{
\abs{p_{xy}^{(n)} - \pi_y} = \abs{P_x\set{X_n = y, n < \tau} - P_{\pi}\set{Y_n = y, n < \tau}} \le 2P_{\delta_x, \pi}(\tau > n). 
}
Since the coupling time is almost surely finite for each initial state $x,y \in \sX$, 
we have $\sum_{n \in \N}P_{\delta_x,\pi}\set{\tau = n} = 1$ and the tail-sum $P_{\delta_x,\pi}\set{\tau > n}$ goes to zero as $n$ grows large, and the result follows. 
\end{itemize}
\end{Proof}
\end{frame}

\begin{frame}{Computing invariant distribution}
When the state space $\sX$ is finite, one can find left eigenvector of probability transition matrix $P$ for the largest eigenvalue~$1$. 
This is the invariant distribution that satisfies the global balance equation $\pi = \pi P$. 
\begin{Definition}<2->
Consider a time homogeneous positive recurrent Markov chain $X:\Omega\to\sX^{\Z_+}$ with probability transition matrix $P$ and invariant distribution $\pi \in \cM(\sX)$. 
For any two disjoint sets $A, B \subseteq \sX$,  
the probability flux from set of nodes $A$ to set of nodes $B$ is defined as $\Phi(A,B) =\sum_{x \in A}\sum_{y \in B}\pi_xp_{xy}$. 
\end{Definition}
\begin{block}{Reminder}<3->
The probability flux from a single node $x$ to single node $y$ is denoted by $\Phi(x,y) = \pi_xp_{xy}$. 
\end{block}
\end{frame}

\begin{frame}
\begin{Definition} 
For a time homogeneous Markov chain $X:\Omega\to\sX^{\Z_+}$ with probability transition matrix $P$ 
represented as the weighted transition graph $G = (\sX, E, w)$, 
a \textbf{cut} is defined as the partition $(\sX_1,\sX_2)$ of nodes.
\end{Definition}
\begin{block}{Lemma}<2->
Probability flux balances across cuts. 
\end{block}
\end{frame}

\begin{frame}
\begin{Proof}
A cut for the state space $\sX$ is given by a partition $(\sX_1,\sX_2)$. 
We show that $\Phi(\sX_1,\sX_2) = \Phi(\sX_2,\sX_1)$. 
To this see, we observe that by exchanging sums from the monotone convergence theorem, and exchanging $x$ and $y$ as running variables, we can write the probability flux $\Phi(\sX_1,\sX_2)$ as 
\begin{align*}
\sum_{y\in \sX_1}\pi_y\sum_{x\notin\sX_1}p_{yx} 
&= \sum_{y\in\sX_1}\pi_y(1- \sum_{x\in\sX_1}p_{yx}) 
= \sum_{y\in\sX_1}\pi_y - \sum_{x\in\sX_1}\sum_{y\in\sX_1}\pi_yp_{yx} \\
&= \sum_{y\in\sX_1}\pi_y - \sum_{x\in\sX_1}(\pi_x- \sum_{y\notin\sX_1}\pi_yp_{yx}) \\
&= \sum_{y\notin \sX_1}\pi_y\sum_{x\in\sX_1}p_{yx}.
\end{align*}
\end{Proof}
\end{frame}

\begin{frame}
\begin{cor}
For any states $y \in \sX$, we have $\pi_y(1-p_{yy}) = \pi_y \sum_{x \neq y}p_{yx} = \sum_{x \neq y}\pi_xp_{xy}$. 
\end{cor}
\begin{Proof}<2->
It follows from probability flux balancing across the cut $(\set{y}, \set{y}^c)$. 
\end{Proof}
\end{frame}

\end{document}
